\documentclass[11pt,a4paper]{article}

\usepackage[margin=1in]{geometry}
\usepackage{xcolor}
\usepackage{listings}
\usepackage{amsmath}
\usepackage{enumitem}
\usepackage[hidelinks]{hyperref}

\definecolor{codebg}{gray}{0.96}
\definecolor{ours}{RGB}{140,70,0}
\definecolor{note}{RGB}{0,90,60}

\lstdefinestyle{json}{%
  basicstyle=\ttfamily\small, backgroundcolor=\color{codebg},
  columns=fullflexible, breaklines=true, showstringspaces=false,
  frame=none, xleftmargin=1em, framexleftmargin=1em,
  aboveskip=0.6em, belowskip=0.6em}

\lstdefinestyle{ml}{%
  language=[Objective]Caml, basicstyle=\ttfamily\small,
  backgroundcolor=\color{codebg}, keywordstyle=\color{blue!60!black},
  commentstyle=\itshape\color{black!45}, columns=fullflexible,
  breaklines=true, showstringspaces=false, frame=none,
  xleftmargin=1em, framexleftmargin=1em, aboveskip=0.6em, belowskip=0.6em}

\lstdefinestyle{lean}{%
  basicstyle=\ttfamily\small, backgroundcolor=\color{codebg},
  columns=fullflexible, breaklines=true, showstringspaces=false,
  frame=none, xleftmargin=1em, framexleftmargin=1em,
  aboveskip=0.6em, belowskip=0.6em}

\newcommand{\fromnote}{\textcolor{note}{\textsf{[design note]}}}
\newcommand{\ourcall}{\textcolor{ours}{\textsf{[our choice]}}}

\setlist[itemize]{itemsep=0.15em, topsep=0.35em}
\setlist[enumerate]{itemsep=0.4em, topsep=0.4em}

\title{\textbf{Tatami --- Phase 1}\\[0.3em]
       \large Shredding JSON into OCaml Modules: the Lean implementation}
\author{Durwasa Chakraborty \and Vimala Soundarapandian \and Vishakh Desai}
\date{Working notes}

\begin{document}
\maketitle

\section{What this document is}

This describes the implementation of Phase~1: the program that reads JSON and
writes OCaml module signatures. The program is written in Lean~4 and compiled
to a binary.

The authority on \emph{what} the generated modules look like is the design
note, \texttt{tatami.pdf}. Its \S2.1 fixes the correspondence between a
relational schema and an OCaml module, and where it speaks we follow it
exactly. It is, however, silent on most of what an implementation must decide:
it says nothing about how a member seen as both an integer and a float should
be typed, nothing about member names that are not legal OCaml identifiers, and
nothing about which inputs to refuse. Everything of that kind is ours.

The two are marked throughout: \fromnote{} for what the note fixes, \ourcall{}
for what we decided.

\paragraph{An earlier version of this document} described a different
approach, in which the shredder was an OCaml program specified by
twenty-two numbered scenarios. That approach was abandoned; the numbering it
introduced (S0--S22, C1--C6) appears nowhere in the current sources and should
not be used. This document replaces it entirely.

\section{Why the implementation is in Lean}

The generator is a Lean program compiled to a native binary. It reads JSON on
standard input or from a file and writes OCaml text.

\paragraph{This is not extraction.} Lean has no OCaml backend, and none is
needed. We are not extracting Lean code into OCaml; we are \emph{generating}
OCaml text, which is ordinary data our program produces. The obstacle that
would apply to extraction does not arise.

\paragraph{The output is a datatype, not a string.} The generator's result
type is \texttt{Ocaml.File} --- a small inductive type expressing exactly the
declarations the design note's signature contains, and nothing else. It cannot
express a function body or any expression at all.

That indirection is the point. Properties of the generated code are statable
about a datatype and are not statable about a string: ``every record field
name is distinct'' is a proposition about \texttt{Ocaml.File}, whereas about a
\texttt{String} it would first require writing an OCaml parser to say anything
at all. Turning the datatype into text is a separate, deliberately dull step.

\paragraph{What stays trusted.} Three things, regardless of what is proved:

\begin{itemize}
  \item the printer, \texttt{File $\to$ String}, which is where proofs stop
        and text begins;
  \item any model we write of OCaml's own rules --- we can check the generated
        code against \emph{our} account of what OCaml accepts, not against the
        compiler;
  \item the reader, \texttt{String $\to$ Doc}, which is the same kind of
        boundary as the printer at the other end.
\end{itemize}

\section{What the design note fixes}

\S2.1, Step~4 gives the correspondence, and we follow all of it:

\begin{center}
\begin{tabular}{@{}ll@{}}
\hline
\textbf{Relational} & \textbf{OCaml} \\
\hline
table & module \\
row & value of type \texttt{t} \\
column & record field \\
primary key & abstract type \texttt{id} \\
foreign key & field holding another module's \texttt{id} \\
joining on it & accessor function \\
\texttt{NOT NULL} & plain field \\
nullable & \texttt{option} field \\
\hline
\end{tabular}
\end{center}

\noindent
Step~3 shows the generated signature for one worked example. Our output for
that input is byte-for-byte identical to the listing in the note:

\begin{lstlisting}[style=json]
{ "a": 1, "b": { "c": 1, "d": 2 } }
\end{lstlisting}

\begin{lstlisting}[style=ml]
(* generated.mli *)

module B : sig
  type id
  type t = { id : id; c : int; d : int }
  val get : id -> t
end

module Root : sig
  type id
  type t = { id : id; a : int; b : B.id }
  val get : id -> t
  val b : t -> B.t
end
\end{lstlisting}

The note also states, without showing it, that ``nested and repeated
substructures are normalized into separate relations linked by generated
keys, with foreign-key constraints materialized rather than left implicit''.
The nested case is the example above. The repeated case is \S\ref{sec:arrays}.

\section{Scalars}

\paragraph{Numbers.} \ourcall{} JSON has one number type; OCaml has
\texttt{int} and \texttt{float}. A member is \texttt{int} only if every value
seen there was written without a decimal point \emph{and} fits OCaml's native
\texttt{int}; otherwise it is \texttt{float}.

\begin{lstlisting}[style=json]
{ "qty": 1 }   { "qty": 2 }        ->  qty : int
{ "p": 1 }     { "p": 2.5 }        ->  p : float      (widened)
{ "x": 1.0 }                       ->  x : float      (written with a point)
{ "id_no": 12345678901234567890 }  ->  id_no : float  (outside native range)
\end{lstlisting}

The distinction is lexical rather than by value, so \texttt{1.0} is a float
even though its value is round. That is recoverable from Lean's parser, which
records a number as a mantissa and an exponent: a decimal point shows up as a
non-zero exponent.

\paragraph{Strings and booleans.} \ourcall{} \texttt{string} and
\texttt{bool}, with nothing to decide.

\paragraph{Reading JSON.} \ourcall{} We read JSON ourselves rather than
using \texttt{Lean.Json}, for two reasons. Its parser stores an object as a
tree map, so a repeated member is discarded before anything can refuse it, and
the order of the rest is lost. It also normalises numbers, so the text a
number was written with cannot be recovered --- and the distinction between
\texttt{1} and \texttt{1.0} rests on exactly that.

Our reader keeps members in order with repeats intact, and keeps each number
as the literal it was written with.

\paragraph{Null and absence.} \ourcall{} Both make a member partial, and the
generated type carries \texttt{option}, which is the note's rule for a
nullable column. The two are not distinguished in the type: a corpus where a
member is sometimes \texttt{null} and one where it is sometimes missing
produce the same module. The counts are kept separately in the report
(\S\ref{sec:report}) precisely because the type cannot carry them.

\paragraph{A member never given a type.} \ourcall{} A member present in every
document but always \texttt{null} has no observed type. It is kept, as
\texttt{unit option} --- the type then records both facts that are true: the
member exists, and nothing is known about it.

\section{Objects}

A nested object becomes a table of its own, named for its path, and the parent
holds a key into it with an accessor that follows the key. \fromnote{} This is
exactly the note's example, generalised to any depth. \ourcall{}

\begin{lstlisting}[style=json]
{ "a": 1, "b": { "c": 1, "inner": { "z": true } } }
\end{lstlisting}

\begin{lstlisting}[style=ml]
module B_inner : sig
  type id
  type t = { id : id; z : bool }
  val get : id -> t
end

module B : sig
  type id
  type t = { id : id; c : int; inner : B_inner.id }
  val get : id -> t
  val inner : t -> B_inner.t
end

module Root : sig ... val b : t -> B.t end
\end{lstlisting}

\paragraph{Module names come from paths.} \ourcall{} A table is identified by
its path from the root, and its module is named for the member names along
that path, mangled and joined: the root is \texttt{Root}, \texttt{.b} is
\texttt{B}, \texttt{.b.inner} is \texttt{B\_inner}.

Distinct tables therefore cannot collide, and the reason is worth recording
because it rests on the mangling: paths are distinct by construction, mangling
is injective, and it escapes \texttt{\_} as \texttt{\_5f}, so a literal
underscore in a member name can never be mistaken for the separator joining
two path segments. Sibling objects using the same member name are unaffected
--- \texttt{\{"x":\{"p":1\},"y":\{"p":"s"\}\}} yields \texttt{X} and
\texttt{Y}.

\paragraph{Ordering.} \ourcall{} A module must be declared before it is
referred to. A child's path is strictly longer than its parent's, so emitting
deepest-first puts every child ahead of the parent that points at it.

\paragraph{An optional nested object} makes both the key and the accessor
optional, following the note's rule for nullable columns:

\begin{lstlisting}[style=ml]
type t = { id : id; a : int; b : B.id option }
val b : t -> B.t option
\end{lstlisting}

\section{Arrays}\label{sec:arrays}

The note calls for repeated substructures to be normalized into separate
relations but shows no example, so the shape below is ours. \ourcall{}

A parent cannot hold a single key to many children, so the direction reverses:
the elements become a table whose rows carry a key back to the parent and an
ordering index, and the parent holds \emph{no column at all} --- only an
accessor returning a list.

\begin{lstlisting}[style=json]
{ "id_no": 1, "items": [ { "sku": "a", "qty": 2 },
                         { "sku": "b", "qty": 1 } ] }
\end{lstlisting}

\begin{lstlisting}[style=ml]
(* generated.mli *)

module rec Items : sig
  type id
  type t = { id : id; parent_id : Root.id; idx : int; qty : int; sku : string }
  val get : id -> t
end

and Root : sig
  type id
  type t = { id : id; id_5fno : int }
  val get : id -> t
  val items : t -> Items.t list
end
\end{lstlisting}

\paragraph{Why \texttt{module rec}.} An array makes the references cyclic:
\texttt{Root} reaches \texttt{Items}, and \texttt{Items} carries a key back to
\texttt{Root}. No ordering of plain module declarations satisfies both, so a
schema containing any array is emitted as a mutually recursive group. A schema
without arrays is emitted as plain modules, which is what keeps the note's own
example byte-identical.

\paragraph{The index is not optional.} JSON array order is significant.
Without \texttt{idx}, \texttt{["a","b"]} and \texttt{["b","a"]} would shred
identically and the order could not be recovered.

\paragraph{Arrays of scalars} take the same shape, with a single
\texttt{value} column in place of the element's members:

\begin{lstlisting}[style=ml]
module rec Tags : sig
  type id
  type t = { id : id; parent_id : Root.id; idx : int; value : string }
  val get : id -> t
end
and Root : sig ... val tags : t -> Tags.t list end
\end{lstlisting}

\paragraph{An array whose elements were never seen} --- empty in every
document --- gives the same table with no \texttt{value} column at all. A null
element makes the value column optional: \texttt{[1, null, 2]} gives
\texttt{value : int option}.

\paragraph{Absent and empty arrays are not distinguished.} \ourcall{} There is
no column in the parent for an array member, so there is nowhere to record
that the member was missing rather than empty. Both give an empty list. This
is a deliberate collapse, and the alternative --- a presence column in the
parent --- would invent a column the note does not sanction.

\section{Markings}

Two shapes cannot be read off the data, and both are supplied by the user
rather than guessed. \ourcall{}

\paragraph{An object whose keys are data} looks exactly like a record. Left
unmarked, \texttt{\{"users": \{"u\_1001": \{...\}, "u\_1002": \{...\}\}\}}
produces one module per user, and a third user adds a third module. Marked, it
produces one table whose rows are keyed by the member name:

\begin{lstlisting}[style=ml]
module rec Users : sig
  type id
  type t = { id : id; parent_id : Root.id; key : string; name : string }
  val get : id -> t
end
and Root : sig ... val users : t -> Users.t list end
\end{lstlisting}

\paragraph{A self-referential shape} looks like a fixed number of nested ones
--- how many depends on how deep the sample went. A comment thread two levels
deep yields three modules; three levels yields four; and a document deeper
than the sample has no table to go in. \ourcall{} This is \emph{not} marked,
and the nesting is taken at face value:

\begin{lstlisting}[style=ml]
module rec Comments_replies_replies : sig ... end
and Comments_replies : sig
  type id
  type t = { id : id; parent_id : Comments.id; idx : int; text : string }
  val get : id -> t
  val replies : t -> Comments_replies_replies.t list
end
and Comments : sig ... val replies : t -> Comments_replies.t list end
and Root : sig ... val comments : t -> Comments.t list end
\end{lstlisting}

A \texttt{recursive} marking did exist, folding a path into the ancestor it
repeats so that both became one table. It was removed. The reason is
\S\ref{sec:gaps}: a table reached from two different places cannot carry a
single \texttt{parent\_id}, and the folded table always is --- rows arrive
both from the collection at the root and from a peer row. Without folding, a
table is identified by its path alone, so its parent and its position column
are functions of that path rather than of the documents, and the question
cannot arise. The cost is the sample-dependence above, and it is recorded as
a known gap rather than settled.

\paragraph{The markings are a file.}

\begin{lstlisting}[style=json]
{ "maps": [".users"] }
\end{lstlisting}

\noindent
passed as \texttt{tatami -{}-config FILE}. Inference then reads
$\mathit{Corpus} \times \mathit{Markings} \rightarrow \mathit{Schema}$: still
deterministic, still a function of structure once the markings are given. A
threshold rule --- treat an object with more than $N$ members as a map --- was
considered and excluded, because member counts grow with the corpus, so it
would make the schema a function of how much data happened to be supplied.

A marking that matches nothing is an error rather than a no-op, so a typo in a
path is reported.

\section{What is refused}

Everything below is our decision; the note does not discuss rejection at all.
\ourcall{}

\begin{center}
\begin{tabular}{@{}p{5.4cm}p{8.2cm}@{}}
\hline
\textbf{Input} & \textbf{Reported as} \\
\hline
A top-level value that is not an object or an array of objects &
  \emph{the top-level value is a number; expected an object, or an array of
  objects.} \\
A member holding two types with no common type &
  \emph{the member .v holds both int and string, which have no common type.} \\
An array with an element that is an array &
  \emph{the array at .m[] has an element that is itself an array. Wrap the
  inner array in an object to give it a name.} \\
An array mixing objects and scalars &
  \emph{the array at .xs[] has both objects and scalars among its elements.} \\
A map entry that is an array &
  \emph{the member at .users\{\} is an array. Map entries may be objects or
  scalars.} \\
A marking that matches nothing &
  \emph{the configuration marks .users as a map, but no object was found
  there.} \\
A repeated member name &
  \emph{the object at .b carries the member c more than once.} \\
A member colliding with a generated column (\texttt{id}, \texttt{parent\_id},
\texttt{idx}) &
  \emph{the member id collides with the generated id column.} \\
An object-valued member named \texttt{get} &
  \emph{... it needs an accessor of that name, which collides with the
  generated get.} \\
Two tables that would take the same module name &
  \emph{two tables would both be called Root.} \\
\hline
\end{tabular}
\end{center}

\noindent
Nested arrays are refused for a structural reason rather than for
convenience: every table is named for the member holding it, and an inner
array sits in element position, which has no member name. There is nothing to
name its table after. Wrapping the inner array in an object supplies the name
and is accepted.

\section{Names}

\ourcall{} A member name is mangled into an OCaml identifier by a total
function: a byte in \texttt{a-z}, \texttt{A-Z}, \texttt{0-9} or \texttt{'}
passes through; every other byte, \emph{including} \texttt{\_}, becomes
\texttt{\_} followed by two lowercase hexadecimal digits; a result not
beginning with a lowercase letter is prefixed \texttt{f\_}; a result that is
an OCaml keyword is suffixed \texttt{\_}.

\begin{lstlisting}[style=json]
my-field  ->  my_2dfield        Name  ->  f_Name
type      ->  type_             1abc  ->  f_1abc
\end{lstlisting}

Escaping \texttt{\_} is what makes the scheme injective, and injectivity is
what stops two distinct members becoming one column. It also makes module
names collision-free, since every mangled name begins with a lowercase letter
and capitalising one character is then a bijection.

Injectivity is currently a property we rely on and have not proved. It is the
first lemma worth doing.

\section{The observation report}\label{sec:report}

Running the binary with \texttt{-{}-json} produces, alongside the module text,
a per-table record of what was seen: how many times an object appeared at that
path, and for each member how many times it held a value, an explicit
\texttt{null}, and how often it was absent.

The denominator differs per table. The root table is visited once per
document; a nested table once per parent that had it; an element table once
per element. A ratio therefore means something different at each depth.

The report is where the \texttt{null}-versus-absent distinction survives,
since the generated \texttt{option} cannot carry it.

\section{Implementation}

\begin{center}
\begin{tabular}{@{}lp{9.3cm}@{}}
\hline
\textbf{File} & \textbf{Holds} \\
\hline
\texttt{Tatami/Doc.lean} & the document type and the JSON reader \\
\texttt{Tatami/Config.lean} & the markings, and reading them from a file \\
\texttt{Tatami/Path.lean} & positions in the document; a path identifies a
  table \\
\texttt{Tatami/Ty.lean} & the types a member can have, and how two
  observations combine \\
\texttt{Tatami/Infer.lean} & reading JSON, walking objects and arrays,
  accumulating per member \\
\texttt{Tatami/Schema.lean} & the handoff: tables and their columns \\
\texttt{Tatami/Mangle.lean} & member names to OCaml identifiers \\
\texttt{Tatami/Ocaml.lean} & the output language \\
\texttt{Tatami/Gen.lean} & the note's correspondence, in code \\
\texttt{Tatami/Print.lean} & output language to text --- the trusted step \\
\texttt{Tatami/Error.lean} & every rejection and its wording \\
\texttt{Main.lean} & the pipeline, the command line, the report \\
\hline
\end{tabular}
\end{center}

\noindent
\texttt{Ty} and \texttt{Infer} know only JSON; \texttt{Ocaml} and
\texttt{Print} know only OCaml; \texttt{Gen} is the only file that touches
both, and \texttt{Schema} is what it passes between them.

\paragraph{The playground.} \texttt{web/playground.html} is a page that shows
the translation as it happens. Served by \texttt{dev/serve.py} it runs the
Lean binary; opened without a server it falls back to an implementation of
the same rules in JavaScript, and says which one produced what is on screen.

Running it both ways is a differential test. All twenty-five cases currently
checked agree, including every rejection.

\paragraph{The output compiles.} Every signature the generator has been run on
has been checked with \texttt{ocamlc -c}, including the mutually recursive
ones. This is a test, not a proof.

\section{What can be guaranteed}

\paragraph{Available now, about the generated code.} That every record field
name within a module is distinct and every referenced type is in scope --- in
short, that the emitted signature always compiles. This is the guarantee worth
having: a duplicate field name is the failure this generator is most likely to
produce, and proving it requires the mangling lemma as a sub-goal.

\paragraph{Available now, about inference.} That the schema does not depend on
the order documents arrive in; that every value observed fits the column it
was assigned; and that the assigned type is the least one that fits. The
lattice facts these rest on (commutativity, idempotence, associativity of the
join) are one-line proofs, since the type has few constructors.

\paragraph{Not available, and why.} Nothing yet can be said about the data
surviving, because \textbf{no data is produced}. The program emits a
description of tables; it never emits a row. Three pieces are missing before a
preservation statement can even be written: something that turns a document
into rows, an interpretation of what \texttt{get} means, and something that
walks rows back into a document.

The design note states in prose what that theorem should say: the original is
no longer stored as a document, and is reconstructed by starting at a
\texttt{Root.t} and calling accessors.

\paragraph{And it would need three exceptions.} Even at present, three things
are lost before a schema is formed: \texttt{null} versus a missing member;
the order of members within an object; and how a number was written. A
round-trip statement can only hold up to those.

\section{Known gaps}\label{sec:gaps}

\begin{enumerate}
  \item \textbf{\texttt{observeObject} and the reader are \texttt{partial}.}
        Lean cannot see that a member of an object is smaller than the object,
        so the recursion is not structural, and the reader's recursion on a
        shrinking input is not visible either. A \texttt{partial} definition
        is opaque to proof, so \texttt{observeObject} must become total before
        anything about inference can be proved --- a size measure on
        \texttt{Doc} and the lemma that members are smaller. This is now
        expressible, which it was not while the parser handed back a tree map.
        The reader is a trusted boundary either way, like the printer.
  \item \textbf{Only one input form.} A single document, or a top-level array
        of them. JSON Lines and concatenated JSON are not read.
  \item \textbf{Self-referential shapes are taken at face value.} There is no
        marking that folds a repeated path into the ancestor it repeats, so a
        comment thread $n$ levels deep yields $n+1$ tables and a document
        deeper than the sample has no table to go in. The schema is therefore
        a function of the sample for such shapes, which is the one place this
        tool does not keep the promise it makes elsewhere.

        The marking is the right design; what it needs first is a parent
        representation able to name more than one source, since a folded
        table's rows arrive both from the collection at the root and from a
        peer row. A variant --- \texttt{Of\_root of Root.id | Of\_comments of
        Comments.id} --- carries that, and carries the position with it, which
        would settle the same question for \texttt{idx} versus \texttt{key}.
        Neither has been built.
\end{enumerate}

\end{document}
